
\documentclass[conference]{IEEEtran}
\IEEEoverridecommandlockouts
\usepackage{cite}
\usepackage{amsmath,amssymb,amsfonts}
\usepackage{graphicx}
\usepackage{booktabs}
\usepackage{array}
\usepackage[spanish]{babel}
\usepackage[utf8]{inputenc}
\usepackage[T1]{fontenc}
\usepackage{url}

\begin{document}

\title{Detección de URLs de \textit{phishing} mediante aprendizaje supervisado sobre el conjunto de datos PhiUSIIL}

\author{
\IEEEauthorblockN{<<Nombre Apellido>>, <<Nombre Apellido>>}
\IEEEauthorblockA{Facultad de Ingeniería, Universidad de Antioquia\\
Curso: Introducción al Machine Learning\\
Correo: <<usuario>>@udea.edu.co}
}

\maketitle

\begin{abstract}
El \textit{phishing} basado en URLs es una de las amenazas de ingeniería social más
extendidas y las listas negras tradicionales no alcanzan a cubrir los sitios creados
cada día. Este trabajo aborda la detección automática de URLs de \textit{phishing} como
un problema de clasificación binaria supervisada, empleando el conjunto de datos público
PhiUSIIL (235.795 URLs y 54 características léxicas, derivadas y del código HTML). En este
avance se presentan la descripción del problema, el análisis exploratorio de los datos y
una revisión del estado del arte sobre el mismo problema y conjunto de datos.
\end{abstract}

\begin{IEEEkeywords}
phishing, detección de URLs, clasificación, aprendizaje supervisado, seguridad informática.
\end{IEEEkeywords}

%=============================================================================
\section{Introducción}
%=============================================================================
El \textit{phishing} es un ataque en el que un sitio web fraudulento suplanta a una entidad
legítima para robar credenciales o información financiera de los usuarios. Las defensas
basadas en listas negras son reactivas: solo bloquean URLs ya reportadas, por lo que fallan
ante los sitios nuevos, que tienen además un tiempo de vida muy corto. Esto motiva el uso de
\textit{machine learning} para aprender, a partir de características de la URL y de la página,
patrones que permitan clasificar un sitio como legítimo o fraudulento de forma automática y
anticipada.

El objetivo de este proyecto es construir y validar un modelo de aprendizaje supervisado que
clasifique una URL como legítima o de \textit{phishing} a partir del conjunto de datos
PhiUSIIL \cite{prasad2024}. Este documento corresponde al avance parcial y cubre dos secciones:
la descripción del problema con su análisis exploratorio (Sección~\ref{sec:problema}) y la
revisión del estado del arte (Sección~\ref{sec:estado}). El código que reproduce los resultados
está disponible en el repositorio del proyecto\footnote{<<URL del repositorio de GitHub>>}.

%=============================================================================
\section{Descripción del problema}
\label{sec:problema}
%=============================================================================

\subsection{Contexto y utilidad de la solución}
Se busca decidir, de forma automática, si una URL corresponde a un sitio legítimo o a uno de
\textit{phishing}. Una solución basada en ML es útil porque generaliza a URLs no vistas (a
diferencia de las listas negras), puede integrarse en navegadores, clientes de correo o
\textit{gateways} de red, y opera en tiempo casi real. El reto es mantener una alta tasa de
detección con una baja tasa de falsos positivos, ya que bloquear sitios legítimos degrada la
experiencia del usuario.

\subsection{Descripción del conjunto de datos}
El conjunto de datos PhiUSIIL \cite{prasad2024}, disponible en el UCI Machine Learning
Repository, contiene \textbf{235.795 muestras} y \textbf{54 características} más la variable
objetivo. Cada muestra representa una URL y su página web asociada. Las clases se distribuyen en
\textbf{134.850 URLs legítimas} (\texttt{label}=1, 57,2\,\%) y \textbf{100.945 URLs de
\textit{phishing}} (\texttt{label}=0, 42,8\,\%), por lo que el problema está
\emph{moderadamente balanceado} (Fig.~\ref{fig:balance}). Las características se agrupan en tres
familias:
\begin{itemize}
  \item \textbf{Léxicas/estructurales de la URL:} \texttt{URLLength}, \texttt{DomainLength},
        \texttt{NoOfSubDomain}, \texttt{NoOfDegitsInURL}, \texttt{SpacialCharRatioInURL},
        \texttt{IsDomainIP}, entre otras.
  \item \textbf{Derivadas:} \texttt{URLSimilarityIndex}, \texttt{CharContinuationRate},
        \texttt{TLDLegitimateProb}, \texttt{URLCharProb}, calculadas por los autores a partir
        de otras variables.
  \item \textbf{Del código HTML:} \texttt{LineOfCode}, \texttt{HasTitle}, \texttt{HasFavicon},
        \texttt{IsResponsive}, \texttt{HasPasswordField}, \texttt{NoOfJS},
        \texttt{HasCopyrightInfo}, \texttt{HasSocialNet}, etc.
\end{itemize}
Además hay cuatro variables de texto: \texttt{URL}, \texttt{Domain}, \texttt{TLD} y
\texttt{Title}. En cuanto a tipos, 50 variables son numéricas (de las cuales 19 son binarias
0/1) y 4 son de texto/categóricas.

\textbf{Datos faltantes.} El conjunto \emph{no presenta valores faltantes} en ninguna columna,
por lo que no se requiere una estrategia de imputación. Tampoco existen filas completamente
duplicadas.

\textbf{Codificación de variables.} Las variables de texto \texttt{URL}, \texttt{Domain} y
\texttt{Title} tienen una cardinalidad cercana al número de muestras (p.\,ej. 235.370 URLs
únicas), por lo que son prácticamente identificadores y no se codifican con \textit{one-hot};
se descartan como variables directas o se reservan para una posible extracción de
características de texto. La variable \texttt{TLD} tiene 695 categorías sin orden; dado que un
\textit{one-hot} generaría cientos de columnas dispersas, se opta por usar la variable numérica
\texttt{TLDLegitimateProb} (que ya codifica el TLD como una probabilidad) o una codificación por
frecuencia. Las variables numéricas tienen escalas muy distintas y colas largas (p.\,ej.
\texttt{LineOfCode} con máximo de 442.666), por lo que se aplicará \textbf{estandarización}
(\texttt{StandardScaler}) para los modelos sensibles a la escala, ajustada solo sobre el
conjunto de entrenamiento.

\begin{figure}[t]
  \centering
  \includegraphics[width=0.85\columnwidth]{figures/01_class_balance.png}
  \caption{Distribución de clases del conjunto PhiUSIIL (57,2\,\% legítimas, 42,8\,\% phishing).}
  \label{fig:balance}
\end{figure}

\subsection{Análisis exploratorio}
El análisis de correlación muestra que varias características del código HTML y derivadas
separan bien las clases: \texttt{URLSimilarityIndex} ($\rho\approx0{,}86$),
\texttt{HasSocialNet} ($0{,}78$), \texttt{HasCopyrightInfo} ($0{,}74$) y \texttt{IsHTTPS}
($0{,}61$). Se identificó una posible \textbf{fuga de información}: todas las URLs con
\texttt{IsHTTPS}=1 están etiquetadas como legítimas, y \texttt{URLSimilarityIndex} está
fuertemente asociada a la etiqueta. Esto podría inflar el desempeño, por lo que en las fases de
modelado se evaluará el efecto de estas variables para garantizar que el modelo generalice. El
detalle del EDA (dimensiones, tipos, duplicados, cardinalidad, distribuciones y correlaciones)
se encuentra en el notebook del repositorio, que comparte las mismas secciones de este informe.

\subsection{Paradigma de aprendizaje}
El problema se enmarca en el \textbf{aprendizaje supervisado}, ya que se dispone de un conjunto
de pares etiquetados $\mathcal{D}=\{(\mathbf{x}_i,y_i)\}_{i=1}^{N}$ con $y_i\in\{0,1\}$. Como la
variable de salida es categórica (legítima / phishing), se trata de un problema de
\textbf{clasificación binaria}. Este paradigma es el apropiado porque el objetivo es predecir
una etiqueta conocida a partir de ejemplos previos, y no descubrir estructura latente (lo que
correspondería a aprendizaje no supervisado). Dado el balance moderado de clases, se usará una
partición \textbf{estratificada} y, además de la exactitud, se reportarán métricas robustas al
desbalance como F1, \textit{balanced accuracy}, MCC y el área bajo la curva ROC (AUC).

%=============================================================================
\section{Estado del arte}
\label{sec:estado}
%=============================================================================
La Tabla~\ref{tab:sota} resume trabajos que abordan la detección de \textit{phishing} por URL
usando el conjunto de datos PhiUSIIL. El trabajo original \cite{prasad2024} introduce el
conjunto de datos y un marco de detección basado en un \emph{índice de similitud} de la URL y
\emph{aprendizaje incremental}, que permite actualizar el modelo ante URLs nuevas.
Posteriormente, \cite{gobinath2026} propone una red neuronal profunda con mecanismo de atención
\textit{multi-head} (ADUIN) orientada a la detección de ataques de día cero, mientras que
\cite{bao2026} compara algoritmos clásicos (Random Forest, SVM, Naïve Bayes y una red MLP) y
reporta exactitudes superiores al 99\,\% sobre PhiUSIIL. Finalmente, \cite{yulianti2026} estudia
métodos de selección de características y muestra que con LASSO ($k=5$ variables) se obtiene un
F1 equivalente al del conjunto completo, lo que apunta a modelos más eficientes sin pérdida de
desempeño.

Una métrica que no se vio en clase y que aparece en estos trabajos es el \textbf{coeficiente
kappa de Cohen}, que mide la concordancia corregida por azar:
\begin{equation}
\kappa = \frac{p_o - p_e}{1 - p_e},
\end{equation}
donde $p_o$ es la proporción de aciertos observada y $p_e$ la proporción de acuerdo esperada por
azar. Las demás métricas reportadas (exactitud, precisión, \textit{recall}, F1 y AUC) se
estudiaron en el curso.

% ---- Tabla comparativa (máx. 1 página para toda la sección) ----
\begin{table*}[t]
\centering
\caption{Trabajos que emplean el conjunto de datos PhiUSIIL.}
\label{tab:sota}
\small
\begin{tabular}{p{2.6cm} p{1.6cm} p{3.0cm} p{2.6cm} p{2.4cm} p{2.6cm}}
\toprule
\textbf{Trabajo} & \textbf{Paradigma} & \textbf{Técnica(s)} & \textbf{Validación} & \textbf{Métricas} & \textbf{Resultados} \\
\midrule
Prasad \& Chandra, 2024 \cite{prasad2024} \newline \textit{Computers \& Security} (Elsevier) &
Supervisado (clasif.) &
Índice de similitud de URL + aprendizaje incremental (MLP) &
Entrenamiento/prueba &
Exactitud, precisión, recall, F1 &
Marco propuesto de alta exactitud \newline (<<confirmar \% del PDF>>) \\
\addlinespace
Gobinath \& Manikandan, 2026 \cite{gobinath2026} \newline \textit{Scientific Reports} (Nature) &
Supervisado (clasif.) &
Red neuronal profunda con atención \textit{multi-head} (ADUIN); \textit{baselines}: SVM, RF, LR &
División estratificada 70/15/15 + 10 repeticiones (t-test, McNemar) &
Exactitud, precisión, recall, F1, FPR, detección día-cero &
95\,\% exactitud; 93\,\% precisión; 92\,\% detección día-cero; FPR 3,5\,\% \\
\addlinespace
Bao \textit{et al.}, 2026 \cite{bao2026} \newline \textit{Computers, Materials \& Continua} &
Supervisado (clasif.) &
Random Forest, SVM, Naïve Bayes (Gaussiano), red MLP &
Hold-out estratificado 80/20 (+ k-fold) &
Exactitud, F1, ROC-AUC, log-loss &
RF 99,99\,\%; MLP 99,98\,\%; SVM 99,97\,\% (exactitud) \\
\addlinespace
Yulianti \& Forenziana, 2026 \cite{yulianti2026} \newline \textit{Formosa J. of Comp. and Inf. Science} &
Supervisado (clasif.) &
Selección de características (filtro, \textit{wrapper}, embebido/LASSO) &
Validación cruzada ($k$-fold) &
F1, desviación estándar por \textit{fold} &
LASSO ($k{=}5$): F1 0,9975 vs. 0,9965 del conjunto completo \\
\bottomrule
\end{tabular}
\end{table*}

\textbf{Observación.} Los desempeños reportados sobre PhiUSIIL son muy altos ($>99\,\%$ en varios
casos), lo que es coherente con la fuerte correlación de algunas características con la etiqueta
observada en el EDA. Esto refuerza la necesidad de un protocolo de validación cuidadoso en las
fases siguientes del proyecto.

%=============================================================================
\begin{thebibliography}{1}
\bibitem{prasad2024}
A. Prasad and S. Chandra, ``PhiUSIIL: A diverse security profile empowered phishing URL
detection framework based on similarity index and incremental learning,''
\emph{Computers \& Security}, vol. 136, p. 103545, 2024,
doi: 10.1016/j.cose.2023.103545.

\bibitem{gobinath2026}
R. Gobinath and S. Manikandan, ``Deep learning-based phishing classification framework for
accurate detection using optimized URL intelligence,'' \emph{Scientific Reports}, vol. 16,
art. 15794, 2026, doi: 10.1038/s41598-026-46481-2.

\bibitem{bao2026}
T. M. Bao, K. Shashvat, N. G. Nhu, and D. Le, ``A comparative analysis of machine learning
algorithms for spam and phishing URL classification,'' \emph{Computers, Materials \& Continua},
vol. 87, no. 2, 2026, doi: 10.32604/cmc.2025.075161.

\bibitem{yulianti2026}
F. Yulianti and J. Forenziana, ``A comparative analysis of feature selection methods for
machine learning-based phishing URL detection,'' \emph{Formosa Journal of Computer and
Information Science}, vol. 5, no. 2, pp. 243--258, 2026, doi: 10.55927/fjcis.v5i2.16783.
\end{thebibliography}

\end{document}
